\begin{minipage}{0.62\linewidth}
    Un générateur de force électromotrice $E=\qty{12}{\V}$ et de résistance interne
    $r=\qty{2}{\ohm}$ est utilisé pour alimenter le <<pont de résistances>>
    représenté sur le schéma ci-dessous.
    \begin{enumerate}
        \item Déterminer la résistance équivalente à l'ensemble
              $R_{\text{éq}}$.~\textbf{(1pt)}
        \item Écrire la loi d'Ohm pour le générateur.~\textbf{(1pt)}
        \item Montrer que l'intensité $I$ du courant débité par le générateur
              s'écrit~: $I=\frac{E}{R_{\text{éq}}+r}$. Calculer sa
              valeur.~\textbf{(1pt)}
    \end{enumerate}

    \begin{donnees}
        $R_{1}=\qty{20}{\ohm}$~; $R_{2}=\qty{5.0}{\ohm}$~;
        $R_{3}=\qty{15}{\ohm}$~; $R_{4}=60 \Omega$.
    \end{donnees}
\end{minipage}
\hfill
\begin{minipage}{0.36\linewidth}
\begin{figure}[H]
    \centering
    \begin{circuitikz}
        \newdimen\d \d=4cm
        \draw (0,0) node[below] {$N$} coordinate (N)
            to[vsource,v_=$(E{,}r)$,*-*] ++(0,\d) node[above] {$P$} coordinate (P)
            to[short,i=$I$] ++(1\d,0)
            to[short,-*] ++(0,-0.1\d) node[above right] {$A$} coordinate(A)
            to[short] ++(-0.2\d,0)
            to[R,l_=$R_{1}$,-*] ++(0,-0.4\d)
            to[R,l_=$R_{4}$] ++(0,-0.4\d)
            to[short,-*] ++(0.2\d,0) node[below right] {$B$} coordinate(B)
            to[short] ++(0,-0.1\d)
            to[short] (N);
        \draw (A)
            to[short] ++(0.2\d,0)
            to[R,l=$R_{2}$,-*] ++(0,-0.4\d)
            to[R,l=$R_{3}$] ++(0,-0.4\d)
            to[short,-*] (B);
    \end{circuitikz}
\end{figure}
\end{minipage}



